% =============================================================================
% ABSTRACT (Tóm tắt nội dung khóa luận) — 150-250 words
% Must include: background, objectives, study design, key findings, conclusion.
% Keywords: 3-6 keywords
% =============================================================================
\chapter*{ABSTRACT}
\phantomsection
\addcontentsline{toc}{chapter}{ABSTRACT}

\textbf{Background.} Despite the proven efficacy of HPV vaccination, substantial disparities in awareness and uptake persist in low- and middle-income countries. In Vietnam, national evidence on the magnitude and determinants of these inequalities among young women remains limited.

\textbf{Objectives.} This study aimed to (1) estimate the weighted prevalence of HPV vaccine awareness and vaccination among Vietnamese women aged 15--29; (2) quantify socioeconomic inequalities using the concentration index; and (3) decompose the observed inequality via the Erreygers method.

\textbf{Methods.} A cross-sectional analysis of the Vietnam Multiple Indicator Cluster Survey 2020--2021 (MICS6) was conducted. The analytic sample comprised 4,102 women aged 15--29. Survey-weighted descriptive statistics, bivariate chi-squared tests, Poisson regression yielding prevalence ratios (PR), the Wagstaff concentration index, and Erreygers decomposition were employed.

\textbf{Results.} Overall, 62.4\% of women had heard of HPV vaccination, while 7.5\% had received the vaccine. The concentration indices for awareness (CI~=~0.149, $p$~<~0.001) and vaccination (CI~=~0.389, $p$~<~0.001) indicated significant pro-rich inequality. Wealth was the largest contributor, explaining 38.8\% and 51.4\% of the inequality in awareness and vaccination, respectively. Education, rural residence, and mass media exposure were substantial contributors.

\textbf{Conclusion.} Marked socioeconomic inequalities characterize HPV vaccine awareness and uptake among young Vietnamese women. Targeted interventions focusing on economically disadvantaged, less educated, rural, and ethnic minority women are urgently needed to advance health equity in cervical cancer prevention.

\textbf{Keywords:} HPV vaccination; socioeconomic inequality; concentration index; Erreygers decomposition; MICS; Vietnam.

\clearpage